\section{Principle VI: Rights and Laws}

I always believed that the privilege of rights also represents the privilege of duties.

The experience of the nature of things established that the domain of common law is very limited and that there is, on the contrary, a multitude of private rights. Each group of men, which has its distinctive life, deserves a distinctive treatment. Whoever ``favors" himself wants to be treated for his own benefit, and when this wish is not granted, the subject suffers, first of all, in the body which he is a part of, and secondarily in himself, due to the disrepute and the slackening which, by all necessity, affects that body.

Right, in order to be imposed and even to subsist, requires that it be exercised, that it be sustained, that it be published. It presupposes activity or it fades, little by little, in the blood and ashes of massacred men and burnt buildings, then in the sublime cold of these empty spaces where the raised voice of the most passionate of the rhetors is going to die out.

That is what the rhetors will never understand. They spend their life personifying ``the right". But men of action and men of analysis ask themselves what this person without subject, this right without a living substrate, can be,: in their turn, by not understanding the antithesis that lingers everywhere, this right that they oppose to force, this force that they want to make the opposite of the right! It makes as much sense to put the triangle and a color in opposition. There are colored triangles, there can be colors spread on triangular surfaces. I do not conceive a right that would be abstract, which would be separated from a moral or material person in whom it exists, and that is to say, from a force.

There is force, more or less strong, which is right; there is force, more or less strong, which is wrong. But a rational being who, without any force, would be the right or would have the right: that is what I cannot conceive.

The right, which needs force in order to be recognized, needs it even more in order to be.

Legal reasoning starts from the principles of the Just and the Unjust, its first notions already represent the second or the third power of a high abstraction, and its definitions, if extensive, are necessarily fluctuating when it is about adapting their generality to practical life: among the multitude of particular facts, often very different, and which sometimes contradict each other, the spirit is more or less inevitably induced to lose sight of the impersonal reasons to stop its choice or even to lead its attention. It is then that, for lack of impersonal reasons, others appear: the personal motive suddenly appears, active and vigilant, and the idea of the right no longer remains clarified and guided except by the idea of the ``\emph{me}", by this ``\emph{me}" which is not without rights, but which claims to have them all and which instinctively steers the processes of thought towards its sole interest, sometimes understood in a tyrannical sense but always, to some degree, unconcerned with a good order, unconsciously favorable to some anarchy.

May the heavens protect me from saying that the Right leads to anarchy, he who wants, on the contrary, to rule it and pacify it! But he is a born warrior. In my opinion, the preference given to the legal method over the empirical method must belong to flourishing societies, strongly based on the principles which erupt from all sides and are obeyed by everyone.

Separate man from his family, his nation, his trade, tell him that he is king, tell him that he is God, and intoxicate him with the idea of Justice, you will see from what heart he will count the wrongs which will be done to him and what will be able to be his leniencies for the wrongs that he will come to do to others! This judge is too partial, a party too interested and too impassioned for it to be reasonable of him to remand theoretically all litigation. The one who recognizes all rights begins by imposing on the entire world all duties, without forgetting the penalties which correspond to all negligence.

That is the true folly of revolutionary individualism, whether it is political, social, or moral. It is impossible that an animal as sensible, as sad, as vulnerable as man, once placed on the interior altar that the dogmatic liberal erects on him, does not believe, nine times out of ten, himself to be the creditor of his fellow men and of the universe, in the place that the most miserable is on the contrary their debtor to the infinite!

This illusion of the creditor on society can only be only encouraged by the absurd metaphysics of Rights.

\paragraph{Clarifications on the Nature of Law} According to a venerable [the Middle Ages] maxim\footnote{From the Middle Ages: \textit{lex consensus populi et constitutione regis fit}.}, the law is made by the action of the sovereign (\textit{constitutione regis}) and by the acceptance of the subject (\textit{consensus populi}). Without doubt our jurists poorly translate these words. They figure that consensus, this must be understood as the agreement of fact, as noisy as the acclamation which follows the coronation of kings, as the simple approving whisper which Homer had follow the word of the leaders, \textit{Epeythymesan axaioi}. But, in the great majority of cases, in the case on which we do not make little stories about, the agreement consists in the fact of raising no important contradiction, of understanding, of carrying out.

To go on to say: “Mister Subject, here is a law that will legally bind you. Do you really want it? Are you sure you want it? We need your signature,” that is exactly to want to inspire him to say “no”, to argue endlessly in order not to agree. The hatred of the new and the spirit of contradiction are so strong in men that we do not put the public good at their mercy.

Yet it remains true that the law must be made to be easily obeyed. A law must be acceptable. The law is not the law when its pronouncement suffices to provoke the people to fear it. It needs a natural and prompt consent.

The party which attacks a law which no one defends can be wrong or right, that is not the question: the author of such a law will not have less committed that political fault of not waiting for the best conditions and most favorable arrangements of the public spirit. Its law is only a decree of a state of siege that it will have to support, arms in hand, without having the right to complain or be surprised.

The spirit of modern laws is far from the spirit of the laws which govern the deeds of real life, and these real deeds, unable to worry about the scraps of paper which distort them or the spider which weaves its web in the legislative mind, continued to develop following the weights, measures, and numbers which compose them.

It is said: “But the deed can be the crime! The deed can be the atrocity! The deed can be the error!”

Definitely: realistic politics and morals do not provide bare facts and beliefs as the types and models of life. But they recommend two points:

First of all, to consider laws depending on which real deeds are chained, for if deeds can be vicious or criminal, the laws of vice and crime are not; the order of causes and effects which preside over realities themselves flawed are not the least flawed in the world, it is even excellent to know and to calculate, knowledge and calculation which alone permit action.

Secondly, action has the chance to be serious and useful only on the condition of aiming for a definite and just goal. It is neither sufficient to have an “ideal” in mind, nor to make a second-rate idea of “right” and “duty”: this moralism, this legalism, this idealism must conform to the ideal, moral, and legal truth. In other words, it is first necessary to be right.

A false idea is a false idea. The will to impose on the world under the pretext that the world must be governed by ideas is an absurd pretence whose application will be forcibly criminal or fatal. Ideas are not equal to each other. All the claimed rights are not valuable. And not all ways of arranging mores are worthy of respect. The modern error comes from that sudden assimilation of contrary systems dreamed by the human mind. It is definitely an easy error for orators and litigants. It allows the latter to support all causes. It lets the former please all audiences. But the people who trust in them pay the price.


\flrightit{Posted on 2012-11-18 by Cologero }

\begin{center}* * *\end{center}

\begin{footnotesize}\begin{sffamily}



\texttt{Avery Morrow on 2012-11-20 at 05:17 said: }

Although it's written as a general principle, I've found that this is an extremely pertinent post in the light of the current Israel-Hamas war (where a lot of claims of one side being ``in the right" are being thrown around), and every time I quote it people get very confused.


\end{sffamily}\end{footnotesize}
